\documentclass[11pt]{letter}

\usepackage[margin=0.5in]{geometry}
\usepackage{hyperref}

\begin{document}

\begin{letter}{}

\opening{Dear Hiring Team,}

I am writing to apply for the Machine Learning Scientist (Graduate/Postdoc Level) - Contract position with Oncoustics. With a Master of Science in Computer Science from the University of Calgary and hands-on experience building Python and PyTorch systems for deep learning, scientific data processing, model evaluation, and research software, I am excited by the opportunity to contribute to machine learning work that can produce real-world health impact.

My graduate research focused on developing machine learning and data-management solutions for large scientific and geospatial datasets. In my digital elevation model super-resolution project, I built Python and PyTorch workflows for data collection, preprocessing, model training, and evaluation. I trained deep ResNet models for 4x reconstruction tasks and assessed model performance using RMSE, slope, TRI, TPI, and wavelet-based metrics. This work required careful debugging, interpretation of model limitations, and practical judgment about whether results generalized beyond convenient random data splits.

I also developed C++ and Python software for scalable processing, storage, and retrieval of spatiotemporal Earth observation data through my work with BigGeo, Mitacs, and the University of Calgary. That research included tensor-based data representations, time-series approximation, multiresolution retrieval, and performance analysis across large raster datasets. Although the domain differs from ultrasound and clinical workflows, the underlying habits are highly relevant to Oncoustics: building reliable pipelines, reasoning about noisy scientific data, evaluating edge cases, and communicating findings clearly with researchers and developers.

The Oncoustics role is especially compelling to me because it combines neural-network design, model debugging, challenging data cases, regression and classification, waveform or time-series signals, and end-to-end ML pipelines. I enjoy the research-engineering loop of turning an idea into runnable code, testing it carefully, understanding failure modes, and improving the system based on evidence. I would bring strong Python, PyTorch, scientific computing, model evaluation, and collaborative research experience, along with genuine motivation to learn the biomedical context quickly.

Thank you for considering my application. I would welcome the opportunity to discuss how my deep learning research background, software engineering experience, and interest in practical healthtech applications can contribute to Oncoustics.

\closing{Sincerely,

Amir Mirzai Golpayegani

\href{mailto:amirgolpaa24@gmail.com}{amirgolpaa24@gmail.com}}

\end{letter}

\end{document}
